\chapter{Introduction}
\label{ch:introduction}
% First prose draft, 2026-10-24. Project evidence reviewed at revision 63ba80d.

Network traffic provides an observable record of how an Internet of Things
(IoT) device communicates with its environment. A device may continue to
exchange syntactically valid packets after its behaviour becomes abnormal:
the change may concern the mixture of protocols, the amount of traffic, or
the services with which it interacts. Detecting such changes requires a
description of expected behaviour and a decision about which departures
deserve attention. This thesis investigates that problem using a compact
statistical representation of each device's traffic.

The system developed in this work, LM-IDNet, models counts of selected packet
categories with a Dirichlet--multinomial (DM) distribution. It combines a
likelihood-based composition score with total-count and Address Resolution
Protocol (ARP) checks, calibrates the decision rule on separate normal
traffic, and supports guarded changes to the fitted model or threshold.
The investigation concerns both the detector and the numerical computation
on which its likelihood scores depend.

\section{Problem and motivation}

The security consequences of compromised IoT devices extend beyond the
devices themselves. Antonakakis et al. documented how the Mirai botnet used
embedded and IoT devices to launch distributed denial-of-service (DDoS) attacks
against Internet services \cite{antonakakis2017}. This case motivates
monitoring device communication even when the device continues to perform
its apparent primary function. An externally observable change in traffic
can provide evidence that warrants investigation without requiring access
to the device's internal software.

Device-specific behaviour offers one possible basis for such monitoring.
Sivanathan et al. showed that network attributes, including activity
patterns and flow characteristics, can help distinguish consumer IoT
devices \cite{sivanathan2019}. These observations support investigating
normal profiles for individual devices. They do not imply that IoT traffic
is constant: user activity, changes in services, and firmware updates can
alter legitimate communication. Behavioural monitoring must therefore
distinguish deviations from a learned profile from evidence that a
particular attack has occurred.

Signature-based detection recognises specified patterns of known activity.
Anomaly detection instead evaluates observations against a model of normal
behaviour \cite{chandola2009}. This can identify departures for which no
attack signature has been supplied, but it also creates a practical
problem: unusual benign activity can produce an alert, while an attack
that resembles normal traffic can remain undetected. Sommer and Paxson
emphasise the difficulties of transferring learned network intrusion
detectors from controlled experiments to operational environments
\cite{sommer2010}. A detector must consequently be assessed through its
false alerts, missed attacks, assumptions, and operating conditions, rather
than through the existence of an anomaly score alone.

A compact representation is attractive when monitoring is intended to run
near the network edge. Local processing is motivated by response-time,
bandwidth, and data-handling considerations \cite{shi2016}. In the present
work, a completed window is represented by four packet counts, so the
statistical decision can be made without retaining packet payloads in the
model. The cost of capturing, validating, and aggregating the source
traffic remains a separate part of the pipeline. A small scoring input
does not by itself establish that the complete system can process a live
network at its packet arrival rate.

Count modelling introduces a further issue. A fixed multinomial model
assumes a common category-probability vector, whereas ordinary traffic can
vary between windows. The DM distribution permits additional variation by
allowing those probabilities to vary around a shared mean profile
\cite{yu2014,alhajbaddar2025analysis}. Its usefulness depends on whether this
variation describes the observed device traffic and whether departures
remain detectable. A model that tolerates ordinary variation may also
tolerate some malicious changes.

Finally, a normal profile may become stale. Updating it from recent
observations can accommodate change, but can also incorporate the activity
that the detector should identify. Cretu et al. examine contaminated
training data for anomaly sensors, while Kloft and Laskov analyse poisoning
of an online centroid detector \cite{cretu2008,kloft2012}. Their work
motivates explicit controls over learning data. It does not establish that
any particular set of controls makes an adaptive DM detector immune to
poisoning.

\section{Thesis objective and scope}

The objective is to design, implement, and evaluate a reproducible
normal-reference anomaly detector for individual IoT devices. The
statistical starting point is the IoT modelling study of Al-Haj Baddar,
Languasco, and Migliardi, which applies DM modelling to D-Link Camera~5
traffic using ten-minute counts of Transmission Control Protocol (TCP),
User Datagram Protocol (UDP), Simple Service Discovery Protocol (SSDP),
and ARP packets \cite{alhajbaddar2025iot}. LM-IDNet adopts this
representation and investigates its use for calibrated anomaly decisions.
The modelling precedent, the parameter estimator, and the numerical
likelihood algorithm are inherited foundations.

The implemented pipeline separates four activities. First, it validates
packet captures, isolates a configured device, and produces time-aligned
count windows. Second, it estimates a normal model from an earlier
presumed-benign period. Third, it uses a distinct benign period to calibrate
the decision rule. Fourth, it scores later windows and records both the
decision and the statistical evidence behind it. Attack annotations are
joined to those decisions only for offline evaluation.

The learning setting is normal-only. Fitting, calibration, and runtime
update admission use no labelled attack examples, but they require traffic
that can serve as a trusted normal reference. The term \emph{unsupervised}
is used in this restricted sense throughout the thesis. In the taxonomy of
Chandola et al., training from normal examples is commonly described as
semi-supervised anomaly detection \cite{chandola2009}. Neither term should
obscure the baseline assumption. Labels are also used to compare completed
experimental runs during development; the absence of labels from the
learning rule is distinct from their absence from the research process.

The detector supports static operation, threshold adaptation, and periodic
parameter refitting. The conservative policy restricts which windows may
enter learning buffers and which candidate changes may become active.
It retains the original detection references during subsequent updates.
The historical, less restrictive policy is retained as an experimental
comparison. Accepted state can be saved, restored, and rolled back, so the
evaluation includes the behaviour of a continuing detector rather than
only an isolated fitting function.

Two dataset roles are kept separate. D-Link captures provide the workload
connected to the inherited numerical study, with the original traces
published by Roy Chowdhury et al. \cite{chowdhury2021data}. Labelled
Chromecast and Samsung camera traces from the UNSW attack-data collection
support detection evaluation; Hamza et al. describe the associated
collection and attack-monitoring study \cite{hamza2019}. The D-Link
captures used here do not provide the attack labels required to estimate
detection precision or recall.

The native Languasco--Migliardi (LM) method evaluates likelihood terms
\cite{languasco2023}; Minka's fixed-point iteration estimates the model
parameters \cite{minka2000}. These roles are examined separately. The
numerical investigation compares LM with the SciPy reference
implementation, while the detection investigation examines the resulting
decisions. Improved numerical evaluation is not assumed to imply improved
attack detection.

Operational measurements cover a supplied-window service in a local
Linux/amd64 Docker environment constrained to one CPU and 512~MiB of
memory. This scope includes processing completed count windows and
attempting permitted updates. It does not include a physical embedded
gateway, an integrated live packet collector, or a claim of lossless
alert delivery. The system emits decisions and alert evidence; automatic
blocking and countermeasure selection are outside the implemented scope.

\section{Research questions}

The study addresses three questions:
\begin{enumerate}
  \item \textbf{RQ1---Detection and representation.} How effectively does a
  device-specific DM model, complemented by volume and ARP checks, identify
  attack-labelled traffic under chronological evaluation, and which
  limitations follow from the count representation and window duration?
  \item \textbf{RQ2---Numerical computation.} How does the native LM
  likelihood backend compare with SciPy in numerical agreement and
  computational cost while preserving the same fitting objective and
  window-score definition?
  \item \textbf{RQ3---Adaptation and safeguards.} How do threshold adaptation
  and periodic refitting affect decisions, and what do learning-buffer
  admission, parameter movement, and candidate rejection reveal about the
  tradeoff between adaptation and retaining a trusted reference?
\end{enumerate}

RQ1 requires confusion counts and precision/recall measures, together with
analysis of attack intervals and the time at which a completed window can
produce an alert. RQ2 requires a common numerical workload and explicit
distinction between model fitting, matrix likelihood evaluation, and
single-window scoring. RQ3 requires the update record as well as the final
detection metrics. Equal static and adaptive predictions cannot establish
a learning benefit when every proposed update is rejected.

These questions do not assume that one policy wins on every device or that
ten minutes is an optimal duration. The duration study considers alternative
aggregation lengths, while the closing evaluation replays the selected
detector from saved starting pairs using fixed settings. Chronological
ordering, development-driven choices, and the independence of an evaluation
are treated as separate matters, in accordance with the experimental-bias
concerns discussed by Pendlebury et al. and Arp et al.
\cite{pendlebury2019,arp2022}.

\section{Contributions}

The contribution is an implemented and evaluated use of established
statistical methods for the stated monitoring problem. It comprises:
\begin{enumerate}
  \item \textbf{A complete count-based detection pipeline.} The work connects
  device filtering and window construction to normal-model fitting,
  independent calibration, and joint composition/volume/ARP decisions.
  Event records retain the observed counts, feature margins, expected
  profile, and model identity needed to inspect a decision.
  \item \textbf{Integration and evaluation of the LM numerical backend.}
  The supplied method is connected to the application's row-wise
  likelihood contract, with validation, sparse-count handling, native
  resource management, and comparison against SciPy. The underlying LM
  approximation and Minka estimator are attributed to their sources.
  \item \textbf{Inspectable adaptive operation.} Static, threshold-adaptive,
  and refitting modes share a scoring path. The conservative policy adds
  explicit admission and movement checks, complete candidate activation,
  persistent accepted state, and rollback. The update records permit
  rejected candidates and contamination risks to be examined alongside
  any accepted changes.
  \item \textbf{A traceable experimental evaluation.} The work compares
  labelled devices, window durations, numerical backends, and adaptation
  policies, and retains the selected-method final-partition outputs.
  The resource study measures the agreed constrained simulation. Negative
  outcomes and the boundaries of each experiment are reported with the
  successful software behaviour.
\end{enumerate}

These contributions do not assert a new probability distribution, a new
fixed-point estimator, or the first use of normal IoT profiles. They also
do not establish general detection performance across arbitrary devices.
Reproducibility is supported by retaining configurations, input and model
fingerprints, predictions, and commands, following the principle that a
reported result should retain its computational provenance
\cite{sandve2013}. The empirical conclusions remain limited by the devices,
attacks, observation periods, and execution environments studied.

\section{Thesis structure}

Chapter~\ref{ch:background} introduces network anomaly detection, categorical
count models, DM estimation, numerical likelihood computation, and related
IoT systems. It identifies the inherited foundations and the engineering
and evaluation questions addressed by LM-IDNet.

Chapter~\ref{ch:implementation} describes the implemented pipeline,
calibration and scoring rules, numerical integration, adaptive policies,
and state lifecycle. Chapter~\ref{ch:experiments} defines the datasets,
chronological partitions, annotation alignment, comparison settings,
metrics, and measurement environments.

Chapter~\ref{ch:results} presents the detection, numerical, adaptation,
window-duration, final-partition, and resource measurements.
Chapter~\ref{ch:discussion} interprets them against the research questions
and examines representation limits, contamination, generalisability, and
deployment assumptions. Chapter~\ref{ch:conclusions} states the supported
conclusions and the remaining research directions. Supporting reproduction
details are collected in Appendix~\ref{app:reproducibility}.
